\chapter{Introduzione}

\section{Contesto e motivazione}

La \textit{Social Network Analysis} (SNA) rappresenta un insieme di metodi e strumenti utili per studiare sistemi complessi descrivibili come reti di entità e relazioni. In tale prospettiva, un fenomeno reale viene modellato tramite un grafo, in cui i nodi rappresentano gli attori del sistema e gli archi descrivono le connessioni esistenti tra essi. Un approccio di questo tipo permette di analizzare la struttura nel suo complesso, individuando i nodi centrali, e possibili gruppi coesi, percorsi strategici e pattern ricorrenti.

La SNA trova applicazione in numerosi contesti, tra cui social media, reti di collaborazione scientifica, reti biologiche, sistemi economici e infrastrutture di trasporto. In questo progetto l’attenzione viene rivolta proprio a quest’ultimo ambito, considerando il trasporto aereo come una rete composta da aeroporti e collegamenti tra di essi. Tale dominio risulta particolarmente adatto a un’analisi di rete, poiché presenta una struttura naturalmente modellabile tramite grafi e permette di osservare dinamiche interessanti legate alla centralità di alcuni nodi, alla presenza di hub e all’emergere di comunità di aeroporti fortemente interconnessi.

\section{Obiettivi del progetto}

L’obiettivo del progetto è, quindi, costruire e analizzare una rete di rotte aeree a partire da un dataset contenente informazioni sui collegamenti tra aeroporti. In particolare, i nodi della rete saranno gli aeroporti, mentre gli archi rappresenteranno le rotte dirette tra un aeroporto di origine e uno di destinazione. Poiché il collegamento tra due aeroporti non è, in generale, simmetrico dal punto di vista della rappresentazione del dato, il fenomeno verrà modellato inizialmente come un grafo diretto. Inoltre, per tener conto della rilevanza di una specifica rotta, il peso di ciascun arco verrà definito sulla base del numero di compagnie aeree distinte che operano quel collegamento.

A partire da questa modellazione, il progetto si propone di affrontare diversi livelli di analisi. In una prima fase verranno studiate le caratteristiche generali della rete, come numero di nodi, numero di archi, densità e distribuzione dei gradi. Successivamente verranno calcolate alcune misure di centralità, con lo scopo di identificare gli aeroporti più importanti sotto differenti punti di vista: numero di connessioni, accessibilità, capacità di intermediazione e rilevanza strutturale. Infine, l’attenzione verrà spostata su aspetti più interni della struttura della rete, attraverso l’individuazione di community, l’analisi del massimo k-core, lo studio di alcune ego-network e una breve osservazione delle clique massimali.

Dal punto di vista metodologico, il lavoro combina una fase iniziale di preprocessing del dataset con una successiva fase di costruzione del grafo e di analisi tramite strumenti di \texttt{Python}. In particolare, verranno utilizzate librerie come \texttt{pandas} per la preparazione dei dati, \texttt{networkx} per la modellazione e l’analisi della rete, e \texttt{matplotlib} per la produzione delle principali visualizzazioni. Per alcune analisi strutturali, come la community detection e il k-core, sarà inoltre impiegata una versione non diretta derivata del grafo, più adatta allo studio di sottostrutture coese.